
\section{Maximum Flow}\label{flow: sec: max flow}\doHeader{Max Flow}

The problem is defined as follows.

\begin{mylist}
\item[\bf input:]
  a \emph{flow network} $(G=(V,E), \textsf{cap}, s, t)$, consisting of

  digraph $G=(V,E)$; edge capacities $\textsf{cap}:E\rightarrow \Rp$;
  \emph{source} $s\in V$ and \emph{sink} $t\in V$

\item[\bf output:]
  a \emph{maximum $s$-$t$ flow} $f$ in $G$ (defined below)
\end{mylist}

Fix any input $(G=(V,E), \textsf{cap}, s, t)$.
What is an $s$-$t$ flow?
Imagine that there is a water source at $s$.
Water must be sent from $s$, through the edges and vertices of the graph, to sink $t$.
Each edge $(u,w)\in E$ can carry at most $\capOf {u,w}$ units of water per second from $u$ to $w$.
For each vertex $v$ other than $s$ or $t$,
all water entering $v$ must leave $v$.
The goal is to maximize the rate at which water is sent.


Formally, an $s$-$t$ flow $f$ in $G$
is a function $f:E\rightarrow \Rp$,
assigning a flow value $f(u,w) \ge 0$ to each edge $(u,w)\in E$,
subject to the following constraints.
For each edge $(u,w)\in E$, the flow $f(u,w)$ on the edge
must be at most the capacity $\capOf {u,w}$.
(These are called the \emph{capacity} constraints.)
For each vertex $v\in V\setminus\{s,t\}$ other than the source or sink,
the flow on the edges into $v$
must be equal to the flow on the edges leaving $v$.
(These are called the \emph{conservation} constraints.)

Any function $f:E\rightarrow \Rp$ meeting these constraints is called a \emph{flow} (for the given flow network).
The \emph{value} (or \emph{size}) of the flow, denoted $|f|$,
is the net flow out of the source.
(As we shall see, this necessarily equals the net flow into the sink.)
The goal is to find a flow $f$ of maximum value.

Here is a restatement of the problem.
Let $\into v = \{ u\in V : (u,v)\in E\}$ denote the in-neighbors of $v$.
Let $\out v = \{ w\in V: (v,w)\in E\}$ denote the out-neighbors.
Using this notation, the problem is as follows:
\[
\begin{LP}{~~\text{maximize~~} |f|
    ~= \sum_{w \in \out s} f(s, w) \,-\, \sum_{u \in \into s} f(u, s)
    \text{~~~subject to:}}
  \big(\forall (u,w)\in E\big) && 0 ~\le~ f(u, w) & ~\le~ \capOf {u,w}
  && \comment{--- capacity constraints} \\[10pt]
  \big(\forall v\in V\setminus\{s,t\}\big)
  && \sum_{w\in \out v} f(v, w)
                                   & ~= \sum_{u\in \into v} f(u, v)
  && \comment{--- conservation constraints}
\end{LP}
\]

\tikzset{
    dot/.style 2 args={fill, circle, inner sep=0pt, label={#1:\scriptsize #2}},
    vertex/.style={circle, draw, inner sep=0pt, minimum size=17pt},
    edge/.style={->,> = latex'}
}

\newcommand{\CAP}[1]{{\small $#1$}}
\newcommand{\FLOW}[2]{$\mathbf{#1}/$\small${#2}$}

Here is one example of a flow network, a flow in it, and a maximum flow in it:
\medskip

\noindent
\centerline{
\begin{tabular}{|c|c|c|}\hline%&&\\
  a flow network &  a flow of value 11  & a maximum flow \\[4pt]
\begin{tikzpicture}[scale=1]
  \node[vertex] (s) at (0,1) {s} ;
  \node[vertex] (a) at (2,2) {a} ;
  \node[vertex] (b) at (2,0) {b} ;
  \node[vertex] (t) at (4,1) {t} ;
  \draw[edge] (s) -- (a) node[midway,above] {\CAP 9}; 
  \draw[edge] (s) -- (b) node[midway,below] {\CAP 5}; 
  \draw[edge] (a) -- (b) node[midway,left] {\CAP 2}; 
  \draw[edge] (a) -- (t) node[midway,above] {\CAP 6}; 
  \draw[edge] (b) -- (t) node[midway,below] {\CAP 8}; 
\end{tikzpicture}
  &
\begin{tikzpicture}[scale=1]
  \node[vertex] (s) at (0,1) {s} ;
  \node[vertex] (a) at (2,2) {a} ;
  \node[vertex] (b) at (2,0) {b} ;
  \node[vertex] (t) at (4,1) {t} ;
  \draw[edge] (s) -- (a) node[midway,above] {\FLOW 6 9}; 
  \draw[edge] (s) -- (b) node[midway,below] {\FLOW 5 5}; 
  \draw[edge] (a) -- (b) node[midway,left] {\FLOW 1 2}; 
  \draw[edge] (a) -- (t) node[midway,above] {\FLOW 5 6}; 
  \draw[edge] (b) -- (t) node[midway,below] {\FLOW 6 8}; 
\end{tikzpicture}
    &
\begin{tikzpicture}[scale=1]
  \node[vertex] (s) at (0,1) {s} ;
  \node[vertex] (a) at (2,2) {a} ;
  \node[vertex] (b) at (2,0) {b} ;
  \node[vertex] (t) at (4,1) {t} ;
  \draw[edge] (s) -- (a) node[midway,above] {\FLOW 8 9}; 
  \draw[edge] (s) -- (b) node[midway,below] {\FLOW 5 5}; 
  \draw[edge] (a) -- (b) node[midway,left] {\FLOW 2 2}; 
  \draw[edge] (a) -- (t) node[midway,above] {\FLOW 6 6}; 
  \draw[edge] (b) -- (t) node[midway,below] {\FLOW 7 8}; 
\end{tikzpicture}
      \\\hline 
\end{tabular}
}

\newpage

\subsection{The Minimum $s$-$t$ Cut problem}

\begin{mylist}
\item[\bf input:]
  a flow network $(G=(V,E), \textsf{cap}, s, t)$

\item[\bf output:]
  an $s$-$t$ \emph{cut} of minimum \emph{capacity} (terms defined below)
\end{mylist}

An \emph{cut} % of the flow network
is a partition of the vertex set $V$ into two parts\footnote
{So $S\subseteq V$ and $T=V\setminus S$.} $(S,T)$.
Given any such cut, define $E(S,T)$ to be the set of edges crossing the cut from $S$ to $T$.
That is, $E(S,T) = \{(u,w) \in E : u\in S, w \in T\} = E\cap (S\times T)$.
Define the \emph{capacity} of the cut, $\capOf {S, T}$,
to be the total capacity of those edges.
For any flow $f$,
define the \emph{flow from $S$ to $T$}, $f(S, T)$,
to be the total flow on those edges.  That is
\[
  \capOf {S, T} = \sum_{(u,w) \in E(S, T)} \capOf {u,w}, \text{~~~and~~~}
  f(S, T) = \sum_{(u,w) \in E(S, T)} f(u,w).
\]
% For example,~$\capOf{\{s\}} = 9 + 5 = 14$,
% $~\capOf{\{b\}} = 8$,
% and $\capOf{\{s,b\}} = 9 + 8 = 17$. 
In the previous example,
\begin{mylist}
\item $\capOf{\{s\}, \{a,b,t\}} = 9 + 5 = 14$,
% $\capOf{\{s,a\}} = 6 + 2 + 5 = 13$. 
\item $\capOf{\{s, b\}, \{a,t\}} = 17$, and
\item $\capOf{\{s,a,b\}, \{t\}} = 6 + 8 = 14$. 
\end{mylist}

An \emph{$s$-$t$ cut} is a cut $(S, T)$ with the additional property that $s\in S$ and $t\in T$.
Next we'll show that the capacity of \emph{any} $s$-$t$ cut
is an upper bound on the value of \emph{any} $s$-$t$ flow.
To prove it we use the following useful property of flows:

\begin{lemma}\label{flow: across cut}
  Let $f$ and $(S, T)$ be any flow and $s$-$t$ cut.
  Then \(|f| = f(S, T) - f(T, S)\).

  That is, the value of the flow equals
  the flow from $S$ to $T$ minus the flow from $T$ to $S$.
\end{lemma}
In the previous example, for the cut $S=\{s,b\}$, the flow $f$ in the third panel
has $f(\{s,b\}, \{a,t\}) - f(\{a,t\}, \{s,b\}) = (8+7) - 2 = 13 = |f|$.

Summing the conservation constraints
over all vertices $v\in S\setminus\{s\}$ gives the equation in the lemma.
You can find the detailed proof at the end of this note.
The bound we mentioned follows:
\begin{lemma}\label{flow: weak}
  Let $f$ and $(S,T)$ be any flow and $s$-$t$ cut.
  Then $|f| \le \capOf {S, T}$.

  That is, the value of any $s$-$t$ flow is at most the capacity of any $s$-$t$ cut.
\end{lemma}
\begin{proof}
  Lemma~\ref{flow: across cut} implies
  $|f| = f(S,T) - f(T,S) \le f(S, T) \le \capOf {S, T}$.

% [review 2026-09-27] fix: the first inequality uses f(T,S) >= 0, not f(S,T) >= 0
  (The two inequalities hold because $f(T, S) \ge 0$,
  and $f$ satisfies the capacity constraints.)
\end{proof}

In the previous example,
the $s$-$t$ cut $S=\{s,a\}$ has capacity 13.
So, by the lemma, every $s$-$t$ flow has value at most 13.
The flow in the third panel achieves value 13,
so must be a maximum flow.

\subsection{A greedy Max-Flow algorithm that doesn't work}
Consider the following greedy algorithm for \prob{Maximum Flow}:

\begin{myframed}
\begin{steps}
  \item[] {\sf greedy-flow}$(G=(V,E), \textsf{cap}, s, t)$

    \step initialize $f[u, w] = 0$ for all $(u,w)\in E$

    \step for any path $p$ in $G$,
    define $\delta(p) = \min\{ \capOf {u,w} - f[u,w] : (u,w) \in p\}$
    \comment{($p$'s residual capacity)}

    \step while there exists an $s$-$t$ path $p$
    with positive residual capacity $\delta(p)>0$:

    \begin{block}
      \step let $p$ be any such path

      \step for each edge $(u,w)\in p$: increase $f[u,w]$ by $\delta(p)$
      \comment{(augment the flow along path $p$)}
    \end{block}

    \step return $f$
  \end{steps}
\end{myframed}

In Step 3, the path $p$ with positive residual capacity is called an \emph{augmenting} path.
Computing one is easy. First delete the \emph{saturated} edges
--- those with $f[u,w] = \capOf{u,w}$.
Then the possible augmenting paths (the $s$-$t$ paths with positive residual capacity)
are the $s$-$t$ paths in the resulting graph.
Such a path can be found (if it exists) in linear time
using depth-first or breadth-first search.

The greedy algorithm maintains the invariant that \emph{$f$ is a flow}
(that is, $f$ satisfies the capacity and conservation constraints).
The invariant holds initially because each edge has flow zero.
In each iteration, in Step 4.2, increasing the flow by $\delta(p)$
along the $s$-$t$ path $p$ maintains the invariant,
because it maintains the capacity constraints and the conservation constraints.
(Verify!)
So, if the algorithm terminates, it returns a flow.
Unfortunately, the algorithm doesn't work ---
it can get stuck before it finds a maximum flow,
as shown in Figure~\ref{flow: fig: greedy}.

\begin{figure}
\centering
\begin{tabular}{|c|c|c|c|}\hline%&&\\
  iteration 1  &  iteration 2  & iteration 3 & iteration 4 \\[4pt]
\begin{tikzpicture}[xscale=1.5]
  \node[vertex] (s) at (0,1) {s} ;
  \node[vertex] (a) at (1,2) {a} ;
  \node[vertex] (b) at (1,0) {b} ;
  \node[vertex] (t) at (2,1) {t} ;
  \draw[edge, ultra thick] (s) -- (a) node[midway,above] {\FLOW 0 9}; 
  \draw[edge] (s) -- (b) node[midway,below] {\FLOW 0 5}; 
  \draw[edge, ultra thick] (a) -- (b) node[midway,left] {\FLOW 0 5}; 
  \draw[edge] (a) -- (t) node[midway,above] {\FLOW 0 6}; 
  \draw[edge, ultra thick] (b) -- (t) node[midway,below] {\FLOW 0 7}; 
\end{tikzpicture}
  &
\begin{tikzpicture}[xscale=1.5]
  \node[vertex] (s) at (0,1) {s} ;
  \node[vertex] (a) at (1,2) {a} ;
  \node[vertex] (b) at (1,0) {b} ;
  \node[vertex] (t) at (2,1) {t} ;
  \draw[edge, ultra thick] (s) -- (a) node[midway,above] {\FLOW 5 9}; 
  \draw[edge] (s) -- (b) node[midway,below] {\FLOW 0 5}; 
  \draw[edge] (a) -- (b) node[midway,left] {\FLOW 5 5}; 
  \draw[edge, ultra thick] (a) -- (t) node[midway,above] {\FLOW 0 6}; 
  \draw[edge] (b) -- (t) node[midway,below] {\FLOW 5 7}; 
\end{tikzpicture}
    &
\begin{tikzpicture}[xscale=1.5]
  \node[vertex] (s) at (0,1) {s} ;
  \node[vertex] (a) at (1,2) {a} ;
  \node[vertex] (b) at (1,0) {b} ;
  \node[vertex] (t) at (2,1) {t} ;
  % [review 2026-09-27] fix: saturated edges are drawn dashed; (s,a) is saturated here too
  \draw[edge, dashed] (s) -- (a) node[midway,above] {\FLOW 9 9}; 
  \draw[edge, ultra thick] (s) -- (b) node[midway,below] {\FLOW 0 5}; 
  \draw[edge, dashed] (a) -- (b) node[midway,left] {\FLOW 5 5}; 
  \draw[edge] (a) -- (t) node[midway,above] {\FLOW 4 6}; 
  \draw[edge, ultra thick] (b) -- (t) node[midway,below] {\FLOW 5 7}; 
\end{tikzpicture}
    &
\begin{tikzpicture}[xscale=1.5]
  \node[vertex] (s) at (0,1) {s} ;
  \node[vertex] (a) at (1,2) {a} ;
  \node[vertex] (b) at (1,0) {b} ;
  \node[vertex] (t) at (2,1) {t} ;
  \draw[edge, dashed] (s) -- (a) node[midway,above] {\FLOW 9 9}; 
  \draw[edge] (s) -- (b) node[midway,below] {\FLOW 2 5}; 
  \draw[edge, dashed] (a) -- (b) node[midway,left] {\FLOW 5 5}; 
  \draw[edge] (a) -- (t) node[midway,above] {\FLOW 4 6}; 
  % [review 2026-09-27] fix: saturated edges are drawn dashed; (b,t) is saturated here, which is why the algorithm is stuck
  \draw[edge, dashed] (b) -- (t) node[midway,below] {\FLOW 7 7}; 
\end{tikzpicture}
      \\\hline 
\end{tabular}
\caption{An example in which the greedy algorithm gets stuck.}\label{flow: fig: greedy}
\end{figure}

\subsection{The Ford-Fulkerson algorithm}
In the remainder of this note, we assume that in $G$ no pair $u,w\in V$ of vertices
has edges in both directions ($|\{(u,w),(w,u)\} \cap E| \le 1$).
This assumption is just to simplify the notation.
It is easy to reduce the general case to this one by ``splitting'' edges as necessary.

Intuitively, in the example, the greedy algorithm failed because it pushed too much flow along edge $(a,b)$.
To fix it, we introduce a mechanism for allowing the algorithm to ``cancel'' previously sent flow.
Namely, 
for every edge $(u,w)\in E$,
if the current flow $f$ sends positive flow $f(u,w) > 0$ on the edge,
we add an artificial \emph{back edge} $(w,u)$, giving it capacity $f(u,w)$.
We also remove all saturated edges.
The resulting graph is the \emph{residual graph} $G_f=(V, E_f)$ for $f$,
where \[E_f = \{(u,w)\in E: f(u,w) < \capOf {u,w} \}
  ~\cup~ \{(u,w) : (w,u) \in E,\, f(w,u) > 0 \}.\]
The edge capacities in $G_f$, called \emph{residual capacities} are
\begin{align*}
  \rescap {u,w} &=
                  \begin{cases}
                    \capOf {u,w} - f(u,w) & \text{ if } (u,w)\in E \text{ and } f(u,w) < \capOf{u,w} \\
                    f(w,u) & \text{ if } (w,u)\in E \text{ and } f(w,u) > 0.
                  \end{cases}
\end{align*}

\newpage

We modify the greedy algorithm to maintain the residual graph $G_f$,
and, in each iteration,
to use any $s$-$t$ path in $G_f$ as the augmenting path.
(Note that all residual capacities are positive.)
This gives the Ford-Fulkerson algorithm, shown in Figure~\ref{flow: fig: ffalg}.

\begin{figure}
\begin{myframed}
\begin{steps}
  \item[] {\sf Ford-Fulkerson}$(G=(V,E), \textsf{cap}, s, t)$

    \step initialize $f[u, w] = 0$ for all $(u,w)\in E$

    \step for any $s$-$t$ path $p$ in $G_f$,
    define $\delta(p) = \min\{ \rescap {u,w} : (u,w) \in p\}$
    \comment{($p$'s \emph{residual capacity})}

    \step while there is an $s$-$t$ path in $G_f$:

    \begin{block}
      \step let $p$ be any such path; let $\Delta=\delta(p)$ \comment{(note $\delta(p) > 0$ by def'n of $G_f$)}

      \step for each edge $(u,w)\in p$:
      \comment{(augment flow along augmenting path $p$ by $\Delta$ units)}
      \step~~~~if $(u,w)\in E$: increase $f[u,w]$ by $\Delta$
      \step~~~~else: decrease $f[w, u]$ by $\Delta$
      \comment{(cancel $\Delta$ units of flow on the reverse edge $(w,u)$)}
    \end{block}

    \step return $f$
  \end{steps}
\end{myframed}
\caption{The Ford-Fulkerson augmenting-paths algorithm for \prob{Maximum Flow}.}\label{flow: fig: ffalg}
\end{figure}

\begin{observation}
  The Ford-Fulkerson algorithm
  maintains the invariant that \emph{$f$ is a flow in $G$.}
\end{observation}
The algorithm maintains the invariant because, in each iteration,
augmenting flow along the path $p$ 
preserves the conservation and capacity constraints.
(To see that this still holds when flow is ``canceled'' on an edge, try some examples.)
Figure~\ref{flow: fig: example} shows the execution of the modified algorithm
on the example from Figure~\ref{flow: fig: greedy}.

\smallskip

\newcommand{\XSCALE}{1.49}
\newcommand{\YSCALE}{0.95}

\begin{figure}
% \noindent
\begin{tabular}{|c||c||c||c|}\hline%&&\\
  % iteration 1  &
                   start of iteration 2  & iteration 3 & iteration 4
                                                & iteration 5 (done)
  \\[4pt]
% \begin{tikzpicture}[yscale=\YSCALE, xscale=\XSCALE]
%   \node[vertex] (s) at (0,1) {s} ;
%   \node[vertex] (a) at (1,2) {a} ;
%   \node[vertex] (b) at (1,0) {b} ;
%   \node[vertex] (t) at (2,1) {t} ;
%   \draw[edge] (s) -- (a) node[midway,above] {\FLOW 0 9}; 
%   \draw[edge] (s) -- (b) node[midway,below] {\FLOW 0 5}; 
%   \draw[edge] (a) -- (b) node[midway,left] {\FLOW 0 5}; 
%   \draw[edge] (a) -- (t) node[midway,above] {\FLOW 0 6}; 
%   \draw[edge] (b) -- (t) node[midway,below] {\FLOW 0 7}; 
% \end{tikzpicture}
%   &  % 2 top
\begin{tikzpicture}[yscale=\YSCALE, xscale=\XSCALE]
  \node[vertex] (s) at (0,1) {s} ; 
  \node[vertex] (a) at (1,2) {a} ;
  \node[vertex] (b) at (1,0) {b} ;
  \node[vertex] (t) at (2,1) {t} ;
  \draw[edge] (s) -- (a) node[midway,above] {\FLOW 5 9}; 
  \draw[edge] (s) -- (b) node[midway,below] {\FLOW 0 5}; 
  \draw[edge] (a) -- (b) node[midway,left] {\FLOW 5 5}; 
  \draw[edge] (a) -- (t) node[midway,above] {\FLOW 0 6}; 
  \draw[edge] (b) -- (t) node[midway,below] {\FLOW 5 7}; 
\end{tikzpicture}
    & % 3 top
\begin{tikzpicture}[yscale=\YSCALE, xscale=\XSCALE]
  \node[vertex] (s) at (0,1) {s} ;
  \node[vertex] (a) at (1,2) {a} ;
  \node[vertex] (b) at (1,0) {b} ;
  \node[vertex] (t) at (2,1) {t} ;
  \draw[edge] (s) -- (a) node[midway,above] {\FLOW 9 9}; 
  \draw[edge] (s) -- (b) node[midway,below] {\FLOW 0 5}; 
  \draw[edge] (a) -- (b) node[midway,left] {\FLOW 5 5}; 
  \draw[edge] (a) -- (t) node[midway,above] {\FLOW 4 6}; 
  \draw[edge] (b) -- (t) node[midway,below] {\FLOW 5 7}; 
\end{tikzpicture}
    & % 4 top
\begin{tikzpicture}[yscale=\YSCALE, xscale=\XSCALE]
  \node[vertex] (s) at (0,1) {s} ;
  \node[vertex] (a) at (1,2) {a} ;
  \node[vertex] (b) at (1,0) {b} ;
  \node[vertex] (t) at (2,1) {t} ;
  \draw[edge] (s) -- (a) node[midway,above] {\FLOW 9 9}; 
  \draw[edge] (s) -- (b) node[midway,below] {\FLOW 2 5}; 
  \draw[edge] (a) -- (b) node[midway,left] {\FLOW 5 5}; 
  \draw[edge] (a) -- (t) node[midway,above] {\FLOW 4 6}; 
  \draw[edge] (b) -- (t) node[midway,below] {\FLOW 7 7}; 
\end{tikzpicture}
    & % 5 top
\begin{tikzpicture}[yscale=\YSCALE, xscale=\XSCALE]
  \node[vertex] (s) at (0,1) {s} ;
  \node[vertex] (a) at (1,2) {a} ;
  \node[vertex] (b) at (1,0) {b} ;
  \node[vertex] (t) at (2,1) {t} ;
  \draw[edge] (s) -- (a) node[midway,above] {\FLOW 9 9}; 
  \draw[edge] (s) -- (b) node[midway,below] {\FLOW 4 5}; 
  \draw[edge] (a) -- (b) node[midway,left] {\FLOW 3 5}; 
  \draw[edge] (a) -- (t) node[midway,above] {\FLOW 6 6}; 
  \draw[edge] (b) -- (t) node[midway,below] {\FLOW 7 7}; 
\end{tikzpicture}
      \\\hline   %%%%%%%%%%%%%%%%%%%%%%%%%%%%%%%% bottom
% \begin{tikzpicture}[yscale=\YSCALE, xscale=\XSCALE]
%   \node[vertex] (s) at (0,1) {s} ;
%   \node[vertex] (a) at (1,2) {a} ;
%   \node[vertex] (b) at (1,0) {b} ;
%   \node[vertex] (t) at (2,1) {t} ;
%   \draw[edge, ultra thick] (s) -- (a) node[midway,above] {\CAP 9}; 
%   \draw[edge] (s) -- (b) node[midway,below] {\CAP 5}; 
%   \draw[edge, ultra thick] (a) -- (b) node[midway,left] {\CAP 5}; 
%   \draw[edge] (a) -- (t) node[midway,above] {\CAP 6}; 
%   \draw[edge, ultra thick] (b) -- (t) node[midway,below] {\CAP 7}; 
% \end{tikzpicture}
%   &  % 2 bottom
\begin{tikzpicture}[yscale=\YSCALE, xscale=\XSCALE]
  \node at (-0.1,2.1) {$G_f$} ; 
  \node[vertex] (s) at (0,1) {s} ;
  \node[vertex] (a) at (1,2) {a} ;
  \node[vertex] (b) at (1,0) {b} ;
  \node[vertex] (t) at (2,1) {t} ;
  \draw[edge, ultra thick] (s) to [bend left=15] node[midway,above] {\CAP 4} (a); 
  \draw[edge] (a) to [bend left=15] node[midway,below] {\CAP 5} (s); 
  \draw[edge] (s) -- (b) node[midway,below] {\CAP 5}; 
  \draw[edge] (b) to [bend right=10] node[midway,left] {\CAP 5} (a); 
  \draw[edge, ultra thick] (a) -- (t) node[midway,above] {\CAP 6}; 
  \draw[edge] (b) to [bend right=10] node[midway,below] {\CAP 2}  (t); 
  \draw[edge] (t) to [bend right=10] node[midway,above] {\CAP 5}  (b); 
\end{tikzpicture}
    & % 3 bottom
\begin{tikzpicture}[yscale=\YSCALE, xscale=\XSCALE]
  \node[vertex] (s) at (0,1) {s} ;
  \node[vertex] (a) at (1,2) {a} ;
  \node[vertex] (b) at (1,0) {b} ;
  \node[vertex] (t) at (2,1) {t} ;
  \draw[edge] (a) to [bend left=10] node[midway,above] {\CAP 9} (s); 
  \draw[edge, ultra thick] (s) -- (b) node[midway,below] {\CAP 5}; 
  \draw[edge] (b) to [bend right=10] node[midway,left] {\CAP 5} (a); 
  \draw[edge] (a) to [bend right=-10] node[midway,above] {\CAP 2} (t); 
  \draw[edge] (t) to [bend right=-10] node[midway,below] {\CAP 4} (a); 
  \draw[edge, ultra thick] (b) to [bend right=10] node[midway,below] {\CAP 2}  (t); 
  \draw[edge] (t) to [bend right=10] node[midway,above] {\CAP 5}  (b); 
\end{tikzpicture}
    & % 4 bottom
\begin{tikzpicture}[yscale=\YSCALE, xscale=\XSCALE]
  \node[vertex] (s) at (0,1) {s} ;
  \node[vertex] (a) at (1,2) {a} ;
  \node[vertex] (b) at (1,0) {b} ;
  \node[vertex] (t) at (2,1) {t} ;
  \draw[edge] (a) to [bend left=10] node[midway,above] {\CAP 9} (s); 
  \draw[edge, ultra thick] (s) to [bend right=10] node[midway,below] {\CAP 3} (b); 
  \draw[edge] (b) to [bend right=10] node[midway,above] {\CAP 2} (s); 
  \draw[edge, ultra thick] (b) to [bend right=10] node[midway,left] {\CAP 5} (a); 
  \draw[edge, ultra thick] (a) to [bend right=-10] node[midway,above] {\CAP 2} (t); 
  \draw[edge] (t) to [bend right=-10] node[midway,below] {\CAP 4} (a); 
  \draw[edge] (t) to [bend right=10] node[midway,below] {\CAP 7}  (b); 
\end{tikzpicture} 
    & % 5 bottom
\begin{tikzpicture}[yscale=\YSCALE, xscale=\XSCALE]
  \begin{scope}
    \clip (-0.30,-0.34) rectangle (2,2.34);
    \draw [rounded corners=13mm, thick, dashed, fill=gray!10,
    opacity=0.8, outer sep=0pt, inner sep=0pt]
    (-0.7,1)--(1.37,2.8)--(1.37,-0.8)--cycle; 
  \end{scope}
  \node[vertex, ultra thick] (s) at (0,1) {s} ;
  \node[vertex, ultra thick] (a) at (1,2) {a} ;
  \node[vertex, ultra thick] (b) at (1,0) {b} ;
  \node[vertex] (t) at (2,1) {t} ;
  \draw[edge] (a) to [bend left=10] node[midway,above] {\CAP 9} (s); 
  \draw[edge] (s) to [bend right=10] node[midway,below] {\CAP 1} (b); 
  \draw[edge] (b) to [bend right=10] node[midway,above] {\CAP 4} (s); 
  \draw[edge] (b) to [bend right=5] node[midway,right] {\CAP 3} (a); 
  \draw[edge] (a) to [bend right=5] node[midway,left] {\CAP 2} (b); 
  \draw[edge] (t) to [bend right=-10] node[midway,below] {\CAP 6} (a); 
  \draw[edge] (t) to [bend right=10] node[midway=0.3,below] {\CAP 7}  (b); 
\end{tikzpicture} 
  \\ \hline 
\end{tabular}
\caption{Example execution of the Ford-Fulkerson algorithm.
  The first row shows the flow $f$ in the graph $G$ at the start of each iteration (2--5).
  The second row shows the residual graph $G_f$ and the augmenting path $p$ in that iteration.
  After Iteration 5 there is no $s$-$t$ path in the residual graph $G_f$
  (as there are no edges leaving the $s$-$t$ cut with $S=\{s, a, b\}$), so the algorithm terminates.
 }\label{flow: fig: example}
\end{figure}

\begin{figure}
\renewcommand{\XSCALE}{1.42}
% \renewcommand{\YSCALE}{1.2}
\centering
\begin{tabular}{|c||c||c|} \hline
  flow network $G$ & final flow $f$ in $G$ & final residual graph $G_f$, and cut \\
\begin{tikzpicture}[yscale=\YSCALE, xscale=\XSCALE]
  \node[vertex] (s) at (0,1) {s} ;
  \node[vertex] (a) at (1,2) {a} ;
  \node[vertex] (b) at (1,0) {b} ; 
  \node[vertex] (c) at (2,2) {c} ; 
  \node[vertex] (d) at (2,0) {d} ; 
  \node[vertex] (t) at (3,1) {t} ;
  \draw[edge] (s) -- (a) node[midway,above] { 9}; 
  \draw[edge] (s) -- (b) node[midway,below] { 1}; 
  \draw[edge] (a) -- (b) node[midway,left] { 5}; 
  \draw[edge] (a) -- (c) node[midway,above] { 2}; 
  \draw[edge] (b) -- (d) node[midway,below] { 3}; 
  \draw[edge] (d) -- (c) node[midway,right] { 2}; 
  \draw[edge] (c) -- (b) node[midway, right] { 1}; 
  \draw[edge] (c) -- (t) node[midway,above] { 6}; 
  \draw[edge] (d) -- (t) node[midway,below] { 1}; 
\end{tikzpicture}
&
\begin{tikzpicture}[yscale=\YSCALE, xscale=\XSCALE]
  \node[vertex] (s) at (0,1) {s} ;
  \node[vertex] (a) at (1,2) {a} ;
  \node[vertex] (b) at (1,0) {b} ; 
  \node[vertex] (c) at (2,2) {c} ; 
  \node[vertex] (d) at (2,0) {d} ; 
  \node[vertex] (t) at (3,1) {t} ;
  \draw[edge] (s) -- (a) node[midway,above] {\FLOW 4 9}; 
  \draw[edge] (s) -- (b) node[midway,below] {\FLOW 1 1}; 
  \draw[edge] (a) -- (b) node[midway,left] {\FLOW 2 5}; 
  \draw[edge] (a) -- (c) node[midway,above] {\FLOW 2 2}; 
  \draw[edge] (b) -- (d) node[midway,below] {\FLOW 3 3}; 
  \draw[edge] (d) -- (c) node[midway,right] {\FLOW 2 2}; 
  \draw[edge] (c) -- (b) node[midway] {\FLOW 0 1}; 
  \draw[edge] (c) -- (t) node[midway,above] {\FLOW 4 6}; 
  \draw[edge] (d) -- (t) node[midway,below] {\FLOW 1 1}; 
\end{tikzpicture}
&
\begin{tikzpicture}[yscale=\YSCALE, xscale=\XSCALE]
  \begin{scope}
    \clip (-0.30,-0.34) rectangle (2,2.34);
    \draw [rounded corners=12mm, thick, dashed, fill=gray!10,
    opacity=0.8, outer sep=0pt, inner sep=0pt]
    (-0.7,1)--(1.4,3)--(1.4,-1)--cycle; 
  \end{scope}
  \node[vertex, ultra thick] (s) at (0,1) {s} ;
  \node[vertex, ultra thick] (a) at (1,2) {a} ;
  \node[vertex, ultra thick] (b) at (1,0) {b} ; 
  \node[vertex] (c) at (2,2) {c} ; 
  \node[vertex] (d) at (2,0) {d} ; 
  \node[vertex] (t) at (3,1) {t} ;
  \draw[edge] (s) to [bend left=10] node[midway,above] { 5} (a); 
  \draw[edge] (a) to [bend left=10] node[midway,below] { 4} (s); 
  \draw[edge] (b) to node[midway,below] { 1} (s); 
  \draw[edge] (a) to [bend left=5] node[midway,right] { 3} (b); 
  \draw[edge] (b) to [bend left=5] node[midway,left] { 2} (a); 
  \draw[edge] (c) to node[midway,below] { 2} (a); 
  \draw[edge] (d) to node[midway,below] { 3} (b); 
  \draw[edge] (c) to node[midway,right] { 2} (d); 
  \draw[edge] (c) -- (b) node[midway, right] {1}; 
  \draw[edge] (c) to [bend left=10] node[midway,above] { 2} (t); 
  \draw[edge] (t) to [bend left=10] node[midway,below] { 4} (c); 
  \draw[edge] (t) -- (d) node[midway,below] {1}; 
\end{tikzpicture} 
  \\ \hline 
\end{tabular}
\caption{An example to illustrate the proof of Theorem~\ref{flow: thm: FF}.}\label{flow: fig: FF}
\end{figure}


\begin{theorem}[Ford \& Fulkerson]\label{flow: thm: FF}
   If the algorithm terminates, it returns a maximum $s$-$t$ flow $f$.
\end{theorem}
\begin{proof}
  Assume the algorithm terminates.  Let $f$ be the flow it returns.
  Let $S$ contain the set of vertices reachable from the source $s$
  in the final residual graph $G_f$.
  Let $T=V\setminus S$.
  The sink $t$ is not in $S$, because at termination
  there is no $s$-$t$ path in the residual graph.
  So $(S,T)$ is an $s$-$t$ cut.
  (In Figure~\ref{flow: fig: FF}, $S=\{s, a, b\}$.)
  
  We use the following observation, which 
  follows from the definition of $S$:
  \emph{in $G_f$ there is no edge $(u,w)$ from $S$ to $T$
    (that is, with $u\in S$ and $w\in T$).}
  (Otherwise, since $u$ is reachable from $s$ in $G_f$,
  $w$ would also be reachable from $s$, so would be in $S$.)
  
  Consider any edge $(u,w)\in E$ from $S$ to $T$ in the original graph $G$.
  By the observation above,
  this edge is not present in the residual graph $G_f$.
  This implies that edge $(u,w)$ must be saturated: $f[u,w] = \capOf {u,w}$.
  Summing over all edges from $S$ to $T$ in $E$, we get $f(S, T) = \capOf {S, T}$.
  (In Figure~\ref{flow: fig: FF}, these edges are $(a, c)$ and $(b, d)$.) 
  
  Consider any edge $(u,w)\in E$ from $T$ to $S$ in the original graph $G$.
  The reverse edge $(w, u)$ (which goes from $S$ to $T$) is not present in the residual graph
  (by the same observation).
  By definition of $G_f$, this implies that edge $(u,w)$ has no flow: $f(u,w) = 0$.
  Summing over all edges from $T$ to $S$ in $E$, we get $f(T, S) = 0$.
  (In Figure~\ref{flow: fig: FF}, there is one such edge, $(c, b)$.) 

  By Lemma~\ref{flow: across cut} and the previous two paragraphs,
  $|f| = f(S,T) - f(T,S) = f(S, T) = \capOf {S, T}$.
  That is, the value of $f$ equals the capacity of the cut.
  By Lemma~\ref{flow: weak},
  no flow has value more than $\capOf {S, T}$.
  So $f$ is a maximum flow.
\end{proof}

\pythonCode{max_flow.py}

\subsection{Consequences: The Max-Flow/Min-Cut Theorem; Integer Flows}
We conclude with two important consequences.
\begin{theorem}[max-flow/min-cut]
  In any flow network, the maximum~$s$-$t$ flow
  value equals the minimum~capacity of any $s$-$t$ cut.
\end{theorem}
\begin{proof}
  Let $v^*$ denote the maximum value of any $s$-$t$ flow.
  (It is well defined, as the set of flows is closed under limit.)
  Let $f$ be a flow of value $v^*$.
  Let $S$ contain the vertices reachable from $s$ in $G_f$.
  % [review 2026-09-27] fix: the proof needs t not in S, which follows from f being a maximum flow (this sentence was missing)
  Then $t\notin S$: otherwise an $s$-$t$ path in $G_f$ would be an augmenting path,
  and augmenting $f$ along it would give a flow of value more than $v^*$.
  Following the reasoning in the proof of Theorem~\ref{flow: thm: FF},
  $(S,T)$ is an $s$-$t$ cut with capacity equal to $|f|$.
  By Lemma~\ref{flow: weak}, every cut has capacity at least $|f| = v^*$.
  So  $\capOf {S,T} = v^*$ is the minimum possible $s$-$t$ cut capacity.
\end{proof}


For many applications, the edge capacities are integers,
and we want an \emph{integer-valued} flow,
where the flow $f(u,w)$ on each edge $(u,w)$ is an integer.
\begin{theorem}
  If the edge capacities are integers, the algorithm 
   returns an integer-valued flow.
\end{theorem}
\begin{proof}
  Assume the edge capacities are integers.
  The algorithm maintains the invariant that $f$ is integer-valued,
  because if $f$ is integer valued and each capacity is an integer,
  then each capacity in the residual graph is a positive integer,
  so $\delta(p)$ is a positive integer in each iteration.
  And the algorithm terminates,
  because in each iteration the value of the flow increases by at least 1,
  but the value never exceeds the capacity of any $s$-$t$ cut.
\end{proof}

Note that this algorithm may take exponentially long.
Indeed, there are flow networks (with irrational edge capacities)
in which the algorithm might not even terminate
(depending on how augmenting paths are chosen)!
But there are variants of the algorithm
that are guaranteed to run in polynomial time.
% [review 2026-09-27] fix: "Dinitz's algorithm" usually names the blocking-flow algorithm; Dinitz discovered the shortest-augmenting-path rule independently
For example, the Edmonds-Karp algorithm (discovered independently by Dinitz)
specializes the Ford-Fulkerson algorithm by choosing the augmenting path in each iteration
to have a minimum number of edges.  (Such a path can be found using breadth-first search.)
With this choice, one can prove there are at most $O(|V|\cdot|E|)$ iterations,
so the running time is $O(|V|\cdot|E|^2)$.

\pythonCode{max_flow.py}

% TODO: fold in the Edmonds-Karp analysis (the lemma that residual distances
% never decrease, and the O(|V||E|)-iterations theorem with its proof) from
% max_flow_Edmonds_Karp.tex, a 2018 CS 5800 note, as a subsection here;
% convert it to the current macros, labels, and \LN references, then delete it.

\subsection{Proof of Lemma~\ref{flow: across cut}}

Fix any flow network $(G=(V,E), \textsf{cap}, s, t)$.

\setcounter{lemma}{0}
\begin{restatement}
\begin{lemma}
  Let $f$ and $(S, T)$ be any flow and $s$-$t$ cut.
  Then \(|f| = f(S, T) - f(T, S)\).

  That is, the value of the flow equals
  the flow from $S$ to $T$ minus the flow from $T$ to $S$.
\end{lemma}
\end{restatement}
\stepcounter{lemma}

\begin{proof} % [Proof of Lemma~\ref{flow: across cut}]
  Fix any such flow $f$ and cost $(S, T)$.
  For ease of notation,
  if, for a pair of vertices $u, w\in V$
  there is no edge $(u, w)\in E$,
  define $f(u, w) = 0$.
  Following this convention,
  for every vertex $v$ other than $s$ or $t$
  the conservation constraint for $v$ is
  \(
    \sum_{w\in V} f(v, w)
    ~= \sum_{u\in V} f(u, v).
  \)
  Then
  \begin{align*}
    |f| & {} = \sum_{w\in V} f(s, w) - \sum_{u\in V} f(u, s) &
    & \comment{(by definition of $|f|$)}
    \\
        & {} = \sum_{v\in S}  \sum_{w\in V} f(v, w)  - \sum_{v\in S}  \sum_{u\in V} f(u, v) &
    & \comment{(adding all conservation constraints for $v\in S\setminus \{s\}$)}
    \\
        & {} = \sum_{v\in S} \sum_{w\in V\setminus S} f(v, w) - \sum_{v\in S} \sum_{u\in V\setminus S} f(u, v) &
    & \comment{(canceling the common terms for edges from $S$ to $S$)}
    \\
        & {} = \sum_{v\in S} \sum_{w\in T} f(v, w)  - \sum_{v\in S} \sum_{u\in T} f(u, v) &
    & \comment{(because $T=V\setminus S$)}
    \\
        & {} = f(S, T) - f(T, S) &
    & \comment{(by definition of $f(S, T)$ and $f(T, S)$).}
  \end{align*}
\end{proof}

\subsection{External sources on \prob{Max Flow}}

\begin{itemize}
\item \DPV~Chapter 7.2.  \CLRS Chapter 26. \KT Chapter 7.  \JE Chapter 10.  MIT lecture video:
  

\URL{https://ocw.mit.edu/courses/electrical-engineering-and-computer-science/6-046j-design-and-analysis-of-algorithms-spring-2015/lecture-videos/lecture-13-incremental-improvement-max-flow-min-cut/}

\end{itemize}

%%% Local Variables:
%%% mode: latex
%%% TeX-master: "flow_reductions_LP"
%%% End:
